\scanpage{228}
point $s_0\in S$ on est ramené encore au cas où $a(s_0)$ est proche de $b(s_0)$, de sorte qu'on peut \deleted{\unsure{donc}} se placer dans un ouvert de $X$ assez petit pour qu'il y ait une connexion. Et on gagne.

On peut aussi noter, \deleted{que} pour prouver la \deleted{formule} nullité de (3) comme conséquence de (7), que l'on a
\[
\deleted{(\unsure{15})}\qquad \sum D_{is}\int_{Z_s}\omega_i=\int_{Z_s}\left(\sum\overline D_i\omega'_i\right)_s
\]
où le deuxième mb a un sens, car $\int_{Z_s}\overline\omega_s=0$ si $\overline\omega_s$ est un bord, et $\left(\sum\overline D_i\omega'_i\right)_s$ est une forme définie modulo un bord. \uline{La formule \unsure{envisagée}} \deleted{équivaut} \inserted{revient à}
\[
(15)\qquad\boxed{D_s\int_{Z_s}\omega'_s=\int_{Z_s}(\overset{\deleted{\unsure{-}}}{D}\omega')_s}
\]
qui donne la propriété transcendante essentielle de la connexion intégrable canonique sur $\underline\Phi$. Nous n'allons pas prouver (15) ici, et l'admettrons, mais notons que
\margin{227}
